\documentclass[12pt,a4paper,UTF8]{ctexart}
\usepackage[a4paper,top=22mm,bottom=22mm,left=22mm,right=22mm,headheight=15pt]{geometry}
\newcommand{\Needspace}[1]{\par\begingroup\dimen0=#1\relax\dimen2=\pagegoal\advance\dimen2 by-\pagetotal\ifdim\dimen0>\dimen2\newpage\fi\endgroup}
\newcommand{\needspace}[1]{\par\begingroup\skip0=\lastskip\vskip-\skip0\dimen0=#1\relax\advance\dimen0 by\skip0\vskip\dimen0\penalty-100\vskip-\dimen0\vskip\skip0\penalty9999\endgroup}
\usepackage{amsmath,amssymb,mathtools,bm}
\usepackage{graphicx,xcolor,booktabs,longtable,array,enumitem,fancyhdr,caption,float,microtype,setspace}
\usepackage[hidelinks,unicode]{hyperref}
\usepackage[most]{tcolorbox}
\definecolor{ContestBlue}{HTML}{17365D}
\hypersetup{pdftitle={第062套},pdfauthor={}}
\setstretch{1.24}
\setlength{\parindent}{2em}
\setlength{\parskip}{0.35em}
\setlength{\emergencystretch}{2em}
\setlist{itemsep=0.35em,topsep=0.4em}
\captionsetup{font=small,labelfont=bf,justification=centering}
\setcounter{secnumdepth}{0}
\pagestyle{fancy}
\fancyhf{}
\fancyhead[L]{\small 第062套}
\fancyfoot[C]{\thepage}
\renewcommand{\headrulewidth}{0.4pt}
\ctexset{
  section={format=\Large\bfseries\color{ContestBlue},beforeskip=1.2em,afterskip=0.7em},
  subsection={format=\large\bfseries\color{ContestBlue},beforeskip=1.0em,afterskip=0.5em},
  subsubsection={format=\normalsize\bfseries,beforeskip=0.8em,afterskip=0.35em}
}
\usepackage{calc}
\newcounter{none}
\let\originalsubsection\subsection
\renewcommand{\subsection}{\needspace{6\baselineskip}\originalsubsection}
\makeatletter
\newcommand{\tightlist}{\setlength{\itemsep}{0pt}\setlength{\parskip}{0pt}}
\newcommand{\maxwidth}{\ifdim\Gin@nat@width>\linewidth\linewidth\else\Gin@nat@width\fi}
\newcommand{\maxheight}{\ifdim\Gin@nat@height>0.72\textheight 0.72\textheight\else\Gin@nat@height\fi}
\newcommand{\pandocbounded}[1]{\begingroup\setkeys{Gin}{width=\maxwidth,height=\maxheight,keepaspectratio}#1\endgroup}
\makeatother
\begin{document}
\begin{titlepage}
  \centering
  \vspace*{58mm}
  {\Huge\bfseries\color{ContestBlue} 第062套\par}
  \vfill
  \vspace*{28mm}
\end{titlepage}
\section{第 1 题}\label{ux7b2c-1-ux9898}

光滑水平桌面上，一小球以速率 \(v_0\)
撞向粗糙库边，速度方向与库边法线成角
\(\theta\)。小球与桌面之间的动摩擦因数为
\(\mu\)，法向碰撞完全弹性，忽略库边形变。

（1）先把小球视为质点，求反弹速度与库边法线的夹角 \(\alpha\)。

（2）再把台球视为半径 \(R\),质量 \(m\)
的均匀实心刚体球。碰撞前质心速率为 \(v_0\)，俯视时以角速度 \(\omega_0\)
逆时针自转；入射角为 \(45^\circ\)，反射角为
\(-30^\circ\)。球在运动过程中不受桌面摩擦，碰撞时忽略库边形变。求库边与球之间的动摩擦因数，并求
\(\omega_0\) 应满足的条件。

\begin{figure}[H]
\centering
\includegraphics[width=0.86\linewidth,height=\textheight,keepaspectratio,alt={图1-1　台球入射,反射与旋转方向的约定}]{../figures/question-01.pdf}
\caption*{图1-1　台球入射,反射与旋转方向的约定}
\end{figure}

\newpage

\section{第 2 题}\label{ux7b2c-2-ux9898}

质量均为 \(m\) 的两个质点 \(A,B\) 固定在长为 \(2l\)
的轻质刚杆两端。杆的中点 \(O\) 系一根原长为 \(l\)
的柔软无弹性轻绳，绳的另一端系质量为 \(m\) 的质点
\(C\)。整个系统置于光滑水平面上。初始时绳松弛，\(C\) 位于距杆 \(l/2\)
处，且以速度 \(v_0\) 平行于刚杆运动；绳将沿图示方向绷直，此时绳与杆成
\(30^\circ\)。

（1）固定杆 \(AB\)，求绳绷直后瞬间 \(C\) 的速度大小和方向。

（2）杆 \(AB\) 自由，求绳绷直后瞬间系统质心的速度大小和方向。

（3）承接（2），分别求绳绷直后瞬间 \(A,B,C\) 的速度大小和方向。

（4）质点之间的碰撞均为完全弹性碰撞。求 \(B,C\)
第一次碰撞后至下一次碰撞前的运动；建议在系统质心参考系中描述。

（5）承接（4），求系统在质心参考系中的运动周期，并画出运动过程。

\begin{figure}[H]
\centering
\includegraphics[width=0.86\linewidth,height=\textheight,keepaspectratio,alt={图2-1　轻杆,松弛轻绳与质点的初始几何}]{../figures/question-02.pdf}
\caption*{图2-1　轻杆,松弛轻绳与质点的初始几何}
\end{figure}

\newpage

\section{第 3 题}\label{ux7b2c-3-ux9898}

如图，均匀六芒星板由一个正六边形和沿其六边向外延伸的六个正三角形组成。把它分成十二个全等小正三角形，每个小正三角形的边长为
\(l\),质量为 \(m\)。六芒星以顶点 \(A\) 与水平地面接触。

（1）求六芒星绕质心且垂直板面的轴的转动惯量 \(J_0\)。

（2）在顶点 \(C\) 处施加水平向右的瞬时冲量 \(I\)。冲击过程中 \(A\)
点可沿地面滑动但不能穿入地面，求冲击结束瞬间六芒星质心的速度和角速度。

（3）改设冲击过程中 \(A\)
点由铰链固定，冲击结束后铰链立即失效。求冲击结束瞬间六芒星质心的速度和角速度。

（4）承接（3）。设地面接触处的摩擦因数为 \(\mu\)，要求六芒星能够绕 \(A\)
点翻越，且 \(A\) 点既不滑动也不脱离地面，求冲量 \(I\) 与 \(\mu\)
应满足的条件。本问不计分。

\begin{figure}[H]
\centering
\includegraphics[width=0.814\linewidth,height=\textheight,keepaspectratio,alt={图3-1　六芒星的几何分割,接触点与冲量位置}]{../figures/question-03.pdf}
\caption*{图3-1　六芒星的几何分割,接触点与冲量位置}
\end{figure}

\newpage

\section{第 4 题}\label{ux7b2c-4-ux9898}

有三根内径相同,粗细均匀的细长通管，长度分别为 \(l_1,l_2,l_3\)，其中
\(l_1<l_2\)。管径足够小，不发生对流；管口尺度仍远大于分子自由程，水蒸气在管内作稳态扩散。左右两端各接一个装有等质量水的容器，系统温度相同；竖直管
\(l_3\) 的上端通向干燥大气。左右支路分别设阀门 \(S_1,S_2\)。

第一次实验关闭 \(S_2\),打开 \(S_1\)，左容器经时间 \(T_1\)
蒸发完。第二次实验关闭 \(S_1\),打开 \(S_2\)，右容器经时间 \(T_2\)
蒸发完。第三次补充相同质量的水后同时打开 \(S_1,S_2\)，系统经时间 \(T_3\)
后两容器均蒸发完。

（1）用 \(l_1,l_2,l_3\) 表示 \(T_2/T_1\)。

（2）用 \(l_1,l_2,l_3\) 表示 \(T_3/T_1\)。

\begin{figure}[H]
\centering
\includegraphics[width=0.84\linewidth,height=\textheight,keepaspectratio,alt={图4-1　三通扩散管,阀门与水容器}]{../figures/question-04.pdf}
\caption*{图4-1　三通扩散管,阀门与水容器}
\end{figure}

\newpage

\section{第 5 题}\label{ux7b2c-5-ux9898}

无穷大接地导体平面上有一个半球形导体凸起，半球半径为
\(R\)，二者彼此导通。上半空间为真空。以球心为原点,竖直向上为极轴建立球坐标系。在
\((a,\theta_0,0)\) 处放置点电荷 \(q\)，其中
\(a>R\)，\(0<\theta_0<\pi/2\)。

（1）以电荷在无穷远处为势能零点，求系统静电能。

（2）在远离电荷分布的场点 \((r,\theta,\phi)\)，其中
\(r\gg R,r\gg a,\theta<\pi/2\)，求电势的远场展开，保留到正比于
\(r^{-3}\) 的项。

（3）若点电荷位置为 \((a,\theta_0,\phi_0(t))\)，其中
\(\phi_0(t)=\omega t\)，求 \(z\) 轴附近远处一点的磁感应强度角向分量
\(B_\phi\)。该点到 \(z\) 轴的距离为 \(\rho\)，满足
\(\rho\ll z\)；暂不考虑电磁扰动传播时间。

（4）若考虑电磁扰动传播的延迟，（3）的结论是否仍成立？简述理由。

\begin{figure}[H]
\centering
\includegraphics[width=0.78\linewidth,height=\textheight,keepaspectratio,alt={图5-1　接地平面,半球凸起与外部点电荷}]{../figures/question-05.pdf}
\caption*{图5-1　接地平面,半球凸起与外部点电荷}
\end{figure}

\newpage

\section{第 6 题}\label{ux7b2c-6-ux9898}

一个均匀带电薄球壳，半径为 \(r\),总电荷为 \(q\)，以角速度 \(\omega\) 绕
\(z\) 轴转动。

（1）求球壳的磁矩 \(\bm m\)。

（2）在匀强磁场 \(\bm B=(0,B_0,0)\) 中，求球壳受到的力矩。

（3）在以速度 \(\bm v=\beta c\,\hat{\bm x}\) 相对原系沿 \(x\)
轴运动的参考系 \(S'\) 中观察，求原来欧拉角为 \((\theta,\phi)\)
的表面元在 \(S'\) 中的面电荷密度 \(\sigma'(\theta,\phi)\)。

（4）求 \(S'\) 中的磁矩 \(\bm m'\) 与电偶极矩 \(\bm p'\)。

（5）求球壳在 \(S'\) 中受到的力与力矩。

提示：电磁场变换为

\[
\begin{gathered}
E_x'=E_x,\qquad
E_y'=\gamma(E_y-vB_z),\qquad
E_z'=\gamma(E_z+vB_y),\\
B_x'=B_x,\qquad
B_y'=\gamma\left(B_y+\frac{v}{c^2}E_z\right),\qquad
B_z'=\gamma\left(B_z-\frac{v}{c^2}E_y\right),
\end{gathered}
\]

其中 \(\gamma=(1-\beta^2)^{-1/2}\)。

\begin{figure}[H]
\centering
\includegraphics[width=0.601\linewidth,height=\textheight,keepaspectratio,alt={图6-1　旋转带电球壳,外磁场与沿 x 方向的参考系变换}]{../figures/question-06.pdf}
\caption*{图6-1　旋转带电球壳,外磁场与沿 \(x\) 方向的参考系变换}
\end{figure}

\newpage

\section{第 7 题}\label{ux7b2c-7-ux9898}

为解释有波纹时湖面上太阳倒影呈狭长椭圆的现象，建立如下小振幅模型。取平静湖面上能把阳光反射入人眼的点为原点，竖直向上为
\(z\) 轴，水平面内指向太阳为 \(y\) 轴。原点附近水面写成

\[
z(x,y,t)=A\cos(k_xx+k_yy-\omega t+\varphi_0),
\]

其中

\[
k_x^2+k_y^2=\left(\frac{2\pi}{\lambda}\right)^2,
\qquad A\ll\lambda.
\]

计算中取
\(A=0.1\,\mathrm m,\lambda=10\,\mathrm m\)。不同位置处波矢方向及局部相位均匀随机，因此每个小区域可视为法向随机的小平面镜。

（1）在 \(t=0\) 时，用 \(k_x,A\) 和局部相位 \(\psi=k_xx+k_yy+\varphi_0\)
表示水面向上单位法向量 \((n_x,n_y,n_z)\) 的 \(n_x\)。

（2）求 \(n_y\)。

（3）求 \(n_z\)。

（4）若法向分布满足

\[
\mathrm dP=f(n_x,n_y)\,\mathrm dn_x\mathrm dn_y,
\]

求 \(f(n_x,n_y)\)。

\needspace{30\baselineskip}

（5）把太阳光视为入射角为 \(\theta\)
的平行光。在人的视野中以平静湖面倒影为原点，水平偏角为
\(\xi\)，向上偏角为 \(\eta\)。对 \(|n_x|,|n_y|\ll1\)
的局部镜面，求反射光视角 \(\xi_0\)。

（6）承接（5），求 \(\eta_0\)。

（7）若视角落在 \([\xi,\xi+\mathrm d\xi],[\eta,\eta+\mathrm d\eta]\)
内的概率为

\[
\mathrm dP=g(\xi,\eta)\,\mathrm d\xi\mathrm d\eta,
\]

求 \(g(\xi,\eta)\)。

\begin{figure}[H]
\centering
\includegraphics[width=0.82\linewidth,height=\textheight,keepaspectratio,alt={图7-1　随机二维波面,局部法向与镜面反射}]{../figures/question-07.pdf}
\caption*{图7-1　随机二维波面,局部法向与镜面反射}
\end{figure}

\newpage

\section{第 8 题}\label{ux7b2c-8-ux9898}

空间中按面心立方结构无限重复排列结点。一个边长为 \(l\)
的常规立方单元中，八个顶点与六个面心均为结点；任意距离为 \(l/\sqrt2\)
的最近邻结点之间都用电阻为 \(r\) 的导线连接。每个结点有十二个最近邻。

（1）求相邻结点 \(A,B\) 之间的等效电阻 \(R_{AB}/r\)。

（2）剪去 \(A,B\) 之间直接连接的电阻，求余下网络在 \(A,B\) 间的等效电阻
\(R_{AB}'/r\)。

（3）记次近邻 \(A,C\) 间电阻为 \(R_{AC}\)，再次近邻 \(A,D\) 间电阻为
\(R_{AD}\)，第四近邻 \(A,E\) 间电阻为 \(R_{AE}\)，求

\[
x=\frac{2R_{AC}+4R_{AD}+R_{AE}}r.
\]

\begin{figure}[H]
\centering
\includegraphics[width=0.82\linewidth,height=\textheight,keepaspectratio,alt={图8-1　面心立方结点,近邻层级与最近邻连线}]{../figures/question-08.pdf}
\caption*{图8-1　面心立方结点,近邻层级与最近邻连线}
\end{figure}
\end{document}
